\chapter{四次元超対称性とWess--Zumino模型}

\lectureoverview{前回の講義の一次元超対称性代数にはハミルトニアンしか現れなかった。ローレンツ不変性は今やエネルギーを空間的運動量と再び結びつけることを要求し、超電荷は二成分スピノルにならなければならない。その代数に到達するために、まず質量ゼロのワイルフェルミオンとそのハミルトニアン\(H=\boldsymbol{\sigma}\cdot\mathbf p\)に戻る。次に超電荷を四次元超空間上の微分演算子として表現し、拘束条件\(\bar D_{\dot\alpha}\Phi=0\)によってカイラル超場を選び出し、超空間積分を用いて具体的な場の理論を構成する。逐次の選択\(\Phi^\dagger\Phi\)、\(\Phi^2\)、\(\Phi^3\)が、伝播、ボソン--フェルミオンの質量の等しさ、湯川相互作用、および関連するスカラー相互作用を生み出す。補助場\(F\)はこれらの関係を透明にし、最後のループ図は超対称性が紫外線感受性を和らげる理由を示す。}

\section{なぜハミルトニアンだけでは不十分なのか}

玩具模型では空間が抑圧され、基本関係は
\begin{equation}
\{Q,Q^\dagger\}=2H.
\label{eq:l8-time-only-algebra}
\end{equation}
であった。この形は時間が唯一の座標である場合には十分である。しかし相対論的な理論では、\(H=P_0\)は四元運動量\(P_\mu\)の一成分である。ローレンツブーストはそれを\(P_i\)と混合する。右辺が\(P_0\)だけを選び出す代数は、共変的に変換することができない。

したがって修正は同時に二つのことを行わなければならない。\(H\)を四つの成分\(P_\mu\)すべてで置き換えなければならず、また左辺に、右辺のローレンツベクトルと対応させることができる添字を備えなければならない。質量のないDirac方程式が必要な手がかりを与える。すなわち、相対論的フェルミオンは二成分スピノルで記述することができ、行列\(\sigma^\mu\)はベクトル添字を一対のスピノル添字に変換する。

\section{質量のないフェルミオン、ヘリシティ、カイラリティ}

質量をもつ粒子と質量のない粒子の物理的区別から始めよう。スピン方向に運動する質量をもつスピン\(\tfrac12\)の粒子を考える。その速さは光速より小さいので、観測者はそれを追い越すことができる。新しい座標系では運動量は反転するがスピンは反転しないため、ヘリシティは反転する。したがって、ローレンツ不変な質量をもつ理論は両方のヘリシティを含まなければならない。

質量のない粒子を追い越すことができる観測者は存在しない。固有ローレンツ変換はその運動量をスピンに対して反転させることができないので、一つの既約な質量のない表現は単一のヘリシティをもつことができる。空間回転は飛行方向を逆にすることはできるが、右手を左手に変えることはない。

\begin{classroomqa}
\audiencequestion{追い越しの議論は、カイラリティではなく本当にヘリシティについてのものですか。\lecturetimestamp{00:05:42}}
\lecturerresponse{はい。ヘリシティは運動量に沿ったスピンの射影であるのに対し、カイラリティはカイラル行列、慣例的に\(\gamma^5\)の固有値である。質量のないモードでは二つのラベルは結び付いているので、この例では言葉が互換的に聞こえるのである。質量をもつ粒子ではそれらは区別される。}
\end{classroomqa}

状態数の数え方には一つの但し書きが必要である。\editorialnote{単一の質量のないWeyl場は、一つの粒子ヘリシティと反対の反粒子ヘリシティをもつ。したがって「状態数が半分」とは、四成分Dirac場から一つの二成分Weyl場への減少を指すのであり、反粒子状態が消えたという意味ではない。}

この冒頭の議論で用いるWeyl場を
\begin{equation}
\chi=
\begin{pmatrix}
\chi_1\\
\chi_2
\end{pmatrix},
\end{equation}
と書く。講義ではこの場を\(\psi_\alpha\)と表記するが、以下でカイラル超場に慣例的に置かれる共役Weyl表現と混同しないよう、ここでは\(\chi\)と呼ぶ。
\begin{equation}
H\chi=\boldsymbol{\sigma}\cdot\mathbf p\,\chi.
\label{eq:l8-weyl-hamiltonian}
\end{equation}
運動量\(\mathbf p\)に対して、行列\(\boldsymbol{\sigma}\cdot\mathbf p\)は固有値\(\pm|\mathbf p|\)をもつ。正振動数の枝が粒子モードを与え、負振動数の枝は下に有界でないスペクトルとして保持されるのではなく、反粒子の構成によって再解釈される。

\begin{figure}[H]
\centering
\includegraphics[width=0.70\textwidth]{lecture_08_figure_02.png}
\lectureframe{00:07:44.271}
\caption{カイラル二成分スピノルとハミルトニアン\(H=\boldsymbol{\sigma}\cdot\mathbf p\)。最後の運動量記号は一部隠れているが、話された式がそれを確定する。}
\end{figure}

左右のラベルもまた慣例に依存する。\editorialnote{黒板では\(\boldsymbol{\sigma}\cdot\mathbf p\)の正の固有値に直接左右を対応させている。どちらのカイラリティがその符号を受け取るかは、\(\sigma^\mu\)、\(\bar\sigma^\mu\)、およびWeyl場の定義に依存する。我々は固有値の記述は保持するが、それに慣例に依存しない左/右のラベルを付けることはしない。}

\(H\)を\(i\partial_t\)に、\(\mathbf p\)を\(-i\nabla\)に置き換えると
\begin{equation}
i\partial_t\chi=-i\boldsymbol{\sigma}\cdot\nabla\chi.
\end{equation}
が得られる。
\begin{equation}
\begin{aligned}
\sigma^\mu&=(\mathbf 1,\boldsymbol{\sigma}),\\
\bar\sigma^\mu&=(\mathbf 1,-\boldsymbol{\sigma}),\\
\eta_{\mu\nu}&=\operatorname{diag}(1,-1,-1,-1).
\end{aligned}
\label{eq:l8-sigma-convention}
\end{equation}
を導入する。
すると、いま得られた方程式は
\begin{equation}
i\sigma^\mu\partial_\mu\chi=0.
\label{eq:l8-weyl-equation}
\end{equation}
となる。共役ワイル表現は\(\bar\sigma^\mu\)を用いる。以下で\(\psi_\alpha\)と名付けるのはこの共役表現であり、その自由方程式は\(i\bar\sigma^\mu\partial_\mu\psi=0\)である。共変性を証明するにはローレンツ群のスピノル表現が必要であるが、ここで重要な事実は構造的なものである。すなわち、4つの行列\(\sigma^\mu\)が1つの時空添字を2つのスピノル添字に結びつける。

\section{四次元超代数}

奇座標は今やワイルスピノルのように変換しなければならない：
\begin{equation}
\theta^\alpha=(\theta^1,\theta^2),
\qquad
\bar\theta^{\dot\alpha}=(\bar\theta^{\dot1},\bar\theta^{\dot2}).
\end{equation}
バーは複素共役を表し、ドットは共役ローレンツ表現を区別する。複素2成分、すなわち実4成分のグラスマン成分が存在する。したがって一般の偶超場
\begin{equation}
\mathcal S=\mathcal S(x^\mu,\theta^\alpha,\bar\theta^{\dot\alpha})
\end{equation}
は有限のグラスマン展開をもつ。その最高次の単項式は
\begin{equation}
\theta^2\bar\theta^2
\equiv
(\theta^1\theta^2)(\bar\theta^{\dot1}\bar\theta^{\dot2}),
\end{equation}
であり、4つの独立な奇座標すべての積である。この単項式の係数は一般に\(D\)成分と呼ばれる。

超電荷もまたスピノル添字をもつ。四次元\(N=1\)超対称性において、それらの代数は
\begin{align}
\{Q_\alpha,\bar Q_{\dot\beta}\}
&=2\sigma^\mu_{\alpha\dot\beta}P_\mu,
\label{eq:l8-superalgebra-main}\\
\{Q_\alpha,Q_\beta\}&=0,
&
\{\bar Q_{\dot\alpha},\bar Q_{\dot\beta}\}&=0,
\label{eq:l8-superalgebra-nilpotent}\\
[P_\mu,Q_\alpha]&=0,
&
[P_\mu,\bar Q_{\dot\alpha}]&=0,
&
[P_\mu,P_\nu]&=0.
\label{eq:l8-superalgebra-commutators}
\end{align}
である。因子2は慣習による。行列
\(\sigma^\mu_{\alpha\dot\beta}\)は\(P_\mu\)のベクトル添字を縮約し、左辺に現れるちょうど2つのスピノル添字を残す。

\begin{classroomqa}
\audiencequestion{シグマ行列の\(\alpha\)と\(\beta\)はその行列添字ですか？\lecturetimestamp{00:22:01}}
\lecturerresponse{はい。より正確には、一方はドットなしワイル添字であり、他方はドット付き共役添字である。\(\mu\)は時空ベクトル添字である。シグマ行列はこれらの記述の間の橋渡しである。}
\end{classroomqa}

整理された記法はその橋渡しを明示する。\editorialnote{講義ではドット付き添字を意図的に省略し、黒板で\(\alpha\beta\)と\(\beta\alpha\)の順序について議論する。式~\eqref{eq:l8-superalgebra-main}と\eqref{eq:l8-superalgebra-commutators}は、なされている議論を変えることなく、標準的なドット付き／ドットなし記法を復元する。}

最後の行は、超電荷が保存され、並進と両立することを述べている。運動量成分が可換であるのは、その微分演算子による表現が可換だからである。以下で用いる微分の慣習では
\begin{equation}
P_\mu=i\partial_\mu,
\end{equation}
であるから、式~\eqref{eq:l8-superalgebra-main}は同様に
\begin{equation}
\{Q_\alpha,\bar Q_{\dot\beta}\}
=2i\sigma^\mu_{\alpha\dot\beta}\partial_\mu.
\end{equation}
と書くこともできる。

\section{超空間微分としての超電荷}

運動量は通常の座標を並進させる。超電荷は奇座標を並進させると同時に時空もずらす。左グラスマン微分を固定し、
\begin{align}
Q_\alpha
&=\partial_\alpha
-i\sigma^\mu_{\alpha\dot\alpha}
\bar\theta^{\dot\alpha}\partial_\mu,
\label{eq:l8-Q}\\
\bar Q_{\dot\alpha}
&=-\bar\partial_{\dot\alpha}
+i\theta^\alpha\sigma^\mu_{\alpha\dot\alpha}\partial_\mu.
\label{eq:l8-Qbar}
\end{align}
と定義する。\(\epsilon^\alpha Q_\alpha+
\bar\epsilon_{\dot\alpha}\bar Q^{\dot\alpha}\)を作用させると、超場の無限小超対称変換が得られる。その成分は独立なグラスマン単項式の係数にほかならず、微分演算子はそれらの係数を混合する。

\begin{classroomqa}
\audiencequestion{省略記法\(\bar\theta\sigma\partial\)において、どの\(x\)微分が意味されているのですか？\lecturetimestamp{00:25:40}}
\lecturerresponse{4つすべてです。\(\bar\theta^{\dot\alpha}\sigma^\mu_{\alpha\dot\alpha}\partial_\mu\)における繰り返された\(\mu\)は和をとられており、したがって\(\sigma^\mu\)の各成分は対応する微分\(\partial/\partial x^\mu\)と対にされている。}
\end{classroomqa}

\begin{classroomqa}
\audiencequestion{共役電荷は別の\(\bar\theta\)ではなく\(\theta\)を含むべきではないですか？\lecturetimestamp{00:26:16}}
\lecturerresponse{はい。\(Q\)は\(\theta\)に関して微分し、\(\bar\theta\)を含む。\(\bar Q\)は\(\bar\theta\)に関して微分し、\(\theta\)を含む。これらは共役な演算子である。}
\end{classroomqa}

\begin{classroomqa}
\audiencequestion{\(\sigma^\mu\)は対称ですか？\lecturetimestamp{00:27:04}}
\lecturerresponse{パウリ行列はエルミートである。\((\mathbf1,\boldsymbol{\sigma})\)の4つの行列のうち3つは対称であり、\(\sigma^2\)は通常の転置の下で反対称である。エルミート性が、ドット付き添字とドットなし添字の配置とともに、関連する共役の性質である。}
\end{classroomqa}

\begin{classroomqa}
\audiencequestion{共役方程式のシグマ行列は添字を逆の順序で担いますか？\lecturetimestamp{00:27:13}}
\lecturerresponse{共役はドット付き表現とドットなし表現を交換するので、2つのラベルを互換可能な行列のスロットとして扱うのではなく、安全な記法は\(\sigma^\mu_{\alpha\dot\alpha}\)とその共役\(\bar\sigma^{\mu\,\dot\alpha\alpha}\)である。}
\end{classroomqa}

式~\eqref{eq:l8-Q}と\eqref{eq:l8-Qbar}の公式は、黒板で用いられた省略記法とは表面的な符号が異なる。それらの符号は、\(P_\mu\)が\(+i\partial_\mu\)であるか\(-i\partial_\mu\)であるか、微分が左から作用するか右から作用するか、そして共役な奇変数がどのように順序づけられるかに関係している。重要なのは、1つの慣習が代数全体を同時に満たすことである。

\section{超対称性を生き延びる拘束条件}

一般的な超場は多くの成分を含む。対称性を壊さずにより小さなクラスを求める。まず通常の類推を示す。ベクトル\(\mathbf V\)に対し、制約
\begin{equation}
V_x=0
\end{equation}
は回転不変ではない。回転の後、\(V_y\)が新しい\(x\)成分に寄与する。対照的に、
\begin{equation}
\mathbf V^2=1
\end{equation}
は不変であり、1つの独立な自由度を除く。有用な制約は、対称性の下で制約された場を制約された場へ写さなければならない。

線形な奇作用素\(\bar D_{\dot\alpha}\)が
\begin{equation}
\bar D_{\dot\alpha}\Phi=0.
\label{eq:l8-chiral-constraint}
\end{equation}
を定義するとする。超対称性変換\(\delta\Phi=Q\Phi\)の後、\(\bar D\)が超電荷と反可換なら同じ条件が成り立つ：
\begin{equation}
\bar D(Q\Phi)=-Q(\bar D\Phi)=0.
\end{equation}
これは、説明なしに現れる定義ではなく、構成の動機となる抽象的整合性テストである。

\begin{classroomqa}
\audiencequestion{最後の式で\(Q\)と\(\bar D\)を入れ替えてよいのはなぜか？\lecturetimestamp{00:36:56}}
\lecturerresponse{それらは反可換になるよう選ばれた奇作用素である。したがって順序を入れ替えると負符号が導入され、\(\bar D\)が\(\Phi\)に達すると元の制約により結果はゼロになる。}
\end{classroomqa}

必要な共変微分は、超電荷に対する時空微分の符号を逆にすることで得られる：
\begin{align}
D_\alpha
&=\partial_\alpha
+i\sigma^\mu_{\alpha\dot\alpha}
\bar\theta^{\dot\alpha}\partial_\mu,
\label{eq:l8-D}\\
\bar D_{\dot\alpha}
&=-\bar\partial_{\dot\alpha}
-i\theta^\alpha\sigma^\mu_{\alpha\dot\alpha}\partial_\mu.
\label{eq:l8-Dbar}
\end{align}
これらは
\begin{align}
\{D_\alpha,\bar D_{\dot\beta}\}
&=-2i\sigma^\mu_{\alpha\dot\beta}\partial_\mu,\\
\{D_\alpha,Q_\beta\}
&=\{D_\alpha,\bar Q_{\dot\beta}\}=0,\\
\{\bar D_{\dot\alpha},Q_\beta\}
&=\{\bar D_{\dot\alpha},\bar Q_{\dot\beta}\}=0.
\end{align}
を満たす。

\begin{classroomqa}
\audiencequestion{\(\bar Q\)から\(\bar D\)への唯一の変化は時空微分項の符号か？\lecturetimestamp{00:37:58}}
\lecturerresponse{そのとおり。その符号反転により、\(D\)をいずれかの超電荷と反可換させると交差項が相殺する。}
\end{classroomqa}

式~\eqref{eq:l8-chiral-constraint}のみを課す。\(D_\alpha\Phi=0\)と\(\bar D_{\dot\alpha}\Phi=0\)の両方を課すと、それらの反可換子を通じて\(\partial_\mu\Phi=0\)が強制され、高々時空に依存しない残部が残る。
\editorialnote{黒板には、両方の制約を課すと\(\Phi=0\)が強制されるとある。代数は厳密に時空独立性を強制し、追加の境界条件または成分条件が場をさらに減らす場合がある。}

\section{カイラル制約の求解}

シフトされた座標を定義する：
\begin{equation}
y^\mu=x^\mu+i\theta\sigma^\mu\bar\theta.
\label{eq:l8-chiral-coordinate}
\end{equation}
上記の左微分規約では、関数
\begin{equation}
\Phi=\Phi(y,\theta)
\label{eq:l8-chiral-solution}
\end{equation}
は\(\bar D_{\dot\alpha}\Phi=0\)を満たす。連鎖律により相殺が明示される。\(y^\mu\)を通して微分すると、\(\bar D_{\dot\alpha}\)に反対符号で現れるのと同じ\(i\theta\sigma^\mu\partial_\mu\Phi\)項が得られる。

したがってカイラル超場は独立な\(\bar\theta\)依存性を持たない。\(\bar\theta\)を含むことはできるが、それは組み合わせ\(y^\mu\)を通してのみである。その有限展開は
\begin{equation}
\Phi(y,\theta)
=\phi(y)+\sqrt2\,\theta\psi(y)+\theta^2F(y).
\label{eq:l8-chiral-expansion-y}
\end{equation}
3つの成分型は複素スカラー\(\phi\)、ワイルフェルミオン\(\psi_\alpha\)、複素スカラー\(F\)である。最後は第2の伝播ボソンではなく補助場である。

\begin{classroomqa}
\audiencequestion{\(x\)はそれに直交する新しい座標\(y\)に置き換えられたのか？\lecturetimestamp{00:49:08}}
\lecturerresponse{いいえ。\(y^\mu\)は式~\eqref{eq:l8-chiral-coordinate}の特定の組み合わせの略記であり、追加の時空方向ではない。\(y\)で作業するとカイラリティが簡単になり、通常の場の理論を構成するときには\(x\)に戻して翻訳する。}
\end{classroomqa}

式~\eqref{eq:l8-chiral-expansion-y}を固定した\(x\)で展開すると
\begin{align}
\Phi(x,\theta,\bar\theta)
={}&\phi
+\sqrt2\,\theta\psi
+\theta^2F
+i\theta\sigma^\mu\bar\theta\,\partial_\mu\phi
\notag\\
&-\frac{i}{\sqrt2}\theta^2
\partial_\mu\psi\,\sigma^\mu\bar\theta
+\frac14\theta^2\bar\theta^2\Box\phi.
\label{eq:l8-chiral-expansion-x}
\end{align}
右辺のすべての成分場は\(x\)で評価される。共役反カイラル場は
\(\bar y^\mu=x^\mu-i\theta\sigma^\mu\bar\theta\)と\(\bar\theta\)に依存する。カイラル超場は一貫して実にはなれない。その共役が反対のカイラリティ制約を満たすからである。

\section{超空間積分}

全超空間作用は
\begin{equation}
S_D=\int d^4x\,d^2\theta\,d^2\bar\theta\;K,
\label{eq:l8-D-term-action}
\end{equation}
の形をとる。ここで\(K\)は実超場である。ベレジン積分はその
\(\theta^2\bar\theta^2\)係数を選び出す。奇積分の後には通常の時空ラグランジアンだけが残る。

第2の不変な構成がある。\(W(\Phi)\)がカイラルなら、
\begin{equation}
S_F=\int d^4x\,d^2\theta\;W(\Phi)
+\int d^4x\,d^2\bar\theta\;\overline{W}(\Phi^\dagger)
\label{eq:l8-F-term-action}
\end{equation}
は超対称である。第1の積分は\(\theta^2\)成分を選び、第2の積分はその共役を選ぶ。

\begin{classroomqa}
\audiencequestion{式が\(\bar\theta\)依存性を持たないとき、\(d^2\theta\)だけで積分してよいのはなぜか？\lecturetimestamp{01:16:30}}
\lecturerresponse{カイラリティは超対称性により保たれ、カイラル超場の最高\(\theta\)成分は時空の全微分だけ変化する。したがってその時空積分は不変であり、全\(D\)成分積分が不変であるのと同じである。}
\end{classroomqa}

\section{自由カイラル多重項}

一見構造のない式
\begin{equation}
S_{\mathrm{free}}
=\int d^4x\,d^2\theta\,d^2\bar\theta\;
\Phi^\dagger\Phi.
\label{eq:l8-free-superaction}
\end{equation}
を取る。2つの因子は反対にシフトされた座標に依存する：
\begin{equation}
\Phi=\Phi(x+i\theta\sigma\bar\theta,\theta),
\qquad
\Phi^\dagger
=\Phi^\dagger(x-i\theta\sigma\bar\theta,\bar\theta).
\end{equation}
時空積分の内部ではダミー変数\(x\)をシフトしてよい。1つのシフトは一方の因子からグラスマン変位を取り除けるが、そのとき他方の相対変位を2倍にする。その冪零変位でのテイラー展開は有限個の項で終わる。

ここが導関数の出所である。超空間の表式
\(\Phi^\dagger\Phi\)
には時空導関数が陽に含まれていないが、それでもずらされた引数が成分場の1階および2階導関数を生み出す。展開を掛け合わせ、最上成分のみを残すと、式~\eqref{eq:l8-sigma-convention}の規約では
\begin{equation}
\mathcal L_{\mathrm{free}}
=\partial_\mu\phi^\ast\partial^\mu\phi
+i\bar\psi\bar\sigma^\mu\partial_\mu\psi
+F^\ast F.
\label{eq:l8-free-component-lagrangian}
\end{equation}
となる。部分積分を除いて、スカラー部分は代わりに概略的に\(-\phi^\ast\Box\phi\)と書くこともでき、これが板書の計算で目指された形である。

\begin{classroomqa}
\audiencequestion{超場のテイラー展開における2行目はどこから来るのですか？\lecturetimestamp{01:03:00}}
\lecturerresponse{それは成分展開\(\Phi(y,\theta)=\phi(y)+\sqrt2\theta\psi(y)+\theta^2F(y)\)から来ます。各成分はその後テイラー展開されます。なぜなら\(y\)は\(x\)と冪零なグラスマン双線形だけ異なるからです。}
\end{classroomqa}

\begin{classroomqa}
\audiencequestion{その最後の成分は\(\theta^2\)を持つべきでしょうか？\lecturetimestamp{01:04:54}}
\lecturerresponse{はい。補助成分は\(F\theta^2\)です。すでに独立な\(\theta\)を両方とも含んでいるため、さらに別の\(\theta\)を伴うテイラー移動は消えます。}
\end{classroomqa}

\begin{figure}[H]
\centering
\includegraphics[width=0.82\textwidth]{lecture_08_figure_03.png}
\lectureframe{01:08:57.900}
\caption{ずらされた座標の展開が、通常のラグランジアンのスカラー項、フェルミオン項、補助項にまとめられる。}
\end{figure}

式~\eqref{eq:l8-free-component-lagrangian}の最初の2つの項は、質量のないクライン--ゴルドン方程式とワイル方程式を与える。最後の項は導関数を含まない：
\begin{equation}
\frac{\partial\mathcal L}{\partial F}=F^\ast=0,
\qquad
\frac{\partial\mathcal L}{\partial F^\ast}=F=0.
\label{eq:l8-free-F-equation}
\end{equation}
それは初期データもオンシェル粒子も導入しない。

\begin{figure}[H]
\centering
\includegraphics[width=0.82\textwidth]{lecture_08_figure_04.png}
\lectureframe{01:11:32.360}
\caption{成分ラグランジアンの最後にある、導関数を含まない補助項\(F^\ast F\)。}
\end{figure}

プロパゲーターの言葉でも同じことが言える。\(F\)に対する逆二次演算子は運動量に依存しない。そのフーリエ変換は局所的である：
\begin{equation}
\langle F(x)F^\ast(y)\rangle
\propto \delta^{(4)}(x-y).
\label{eq:l8-F-contact}
\end{equation}
\(F\)線は図を組織するかもしれないが、その端点は一致する。それは接触収縮であり、物理的粒子の軌跡ではない。

\section{二次超ポテンシャルと等しい質量}

伝播は\(D\)項から来た。質量はカイラル積分から来る。従来の規格化
\begin{equation}
W_2(\Phi)=\frac12m\Phi^2.
\end{equation}
を選ぶ。\(\theta^2\)成分は
\begin{equation}
\left.W_2(\Phi)\right|_{\theta^2}
=m\phi F-\frac12m\psi\psi.
\end{equation}
であり、共役項と合わせると、成分ラグランジアンは
\begin{align}
\mathcal L
={}&
\partial_\mu\phi^\ast\partial^\mu\phi
+i\bar\psi\bar\sigma^\mu\partial_\mu\psi
+F^\ast F
\notag\\
&+m\phi F+m^\ast\phi^\ast F^\ast
-\frac12m\psi\psi
-\frac12m^\ast\bar\psi\bar\psi.
\label{eq:l8-mass-offshell}
\end{align}
となる。フェルミオン双線形はワイル場のマヨラナ質量である。\(F\)方程式は代数的である：
\begin{equation}
F^\ast=-m\phi,
\qquad
F=-m^\ast\phi^\ast.
\label{eq:l8-mass-F-equations}
\end{equation}
代入すると
\begin{equation}
\begin{aligned}
\mathcal L_{\mathrm{on\ shell}}={}&
\partial_\mu\phi^\ast\partial^\mu\phi
+i\bar\psi\bar\sigma^\mu\partial_\mu\psi\\
&-|m|^2|\phi|^2\\
&-\frac12\left(m\psi\psi+m^\ast\bar\psi\bar\psi\right).
\end{aligned}
\label{eq:l8-mass-onshell}
\end{equation}
が得られる。したがってスカラー質量とフェルミオン質量はともに\(|m|\)である。板書の図式的な言葉では、スカラーは1つの\(m\phi F\)挿入を通じて局所的に\(F\)に変わり、直ちにその共役を通じて戻る。2つの挿入が\(|m|^2|\phi|^2\)を与え、これは式~\eqref{eq:l8-mass-F-equations}を解いて見出されたスカラー質量項に他ならない。

規格化を今や正確に述べることができる。\editorialnote{講義では超場を掛け合わせる際に数値係数を意図的に省略している。ここで用いられた因子\(1/2\)、\(1/3\)、\(\sqrt2\)は従来のウェス--ズミノ規格化であり、それらは板書の導出のあらゆる定性的な段階を保ちつつ、成分の質量と結合を正確にする。}

\section{三次超ポテンシャルと相互作用}

次に
\begin{equation}
W(\Phi)=\frac12m\Phi^2+\frac13g\Phi^3.
\label{eq:l8-superpotential}
\end{equation}
を加える。一般の正則な\(W\)に対して、\(F\)成分は
\begin{equation}
\left.W(\Phi)\right|_{\theta^2}
=F W'(\phi)-\frac12W''(\phi)\psi\psi.
\label{eq:l8-W-components}
\end{equation}
である。ここで
\begin{equation}
W'(\phi)=m\phi+g\phi^2,
\qquad
W''(\phi)=m+2g\phi.
\end{equation}
である。\(F\)を消去する前、三次項は
\begin{equation}
g\phi^2F-g\phi\psi\psi+\text{h.c.}
\label{eq:l8-cubic-offshell}
\end{equation}
を含む。第2項は1つのスカラーを2本のフェルミオン線に結びつける湯川頂点である。第1項は2つのスカラーを補助線に結びつける。

\begin{figure}[H]
\centering
\begin{tikzpicture}[scale=0.95,>=Latex]
  \coordinate (vone) at (0,0);
  \draw[->,thick] (-2,1.1) -- (vone);
  \draw[->,thick] (vone) -- (-2,-1.1);
  \draw[thick,dashed] (vone) -- (2,0);
  \node[left] at (-2,1.1) {\(\psi\)};
  \node[left] at (-2,-1.1) {\(\psi\)};
  \node[right] at (2,0) {\(\phi\)};
  \node[below=8pt of vone] {\(-g\phi\psi\psi\)};
\end{tikzpicture}
\caption{湯川相互作用の教育的な描き直し：2本のフェルミオン脚と1本のスカラー脚が、\(g\)に比例する結合で出会う。}
\end{figure}

完全な補助場の方程式は
\begin{equation}
F^\ast=-W'(\phi)=-(m\phi+g\phi^2).
\label{eq:l8-general-F-equation}
\end{equation}
である。これを消去するとスカラーポテンシャル
\begin{equation}
V(\phi)=|W'(\phi)|^2
=|m\phi+g\phi^2|^2.
\label{eq:l8-scalar-potential}
\end{equation}
が生じる。質量項のほかに、これは\(m^\ast g\)に比例する三次スカラー項と、四次項
\begin{equation}
V(\phi)\supset |g|^2|\phi|^4.
\label{eq:l8-quartic}
\end{equation}
を含む。補助線の描像では、2つの\(\phi^2F\)頂点が局所的な\(F\)収縮で結ばれる。式~\eqref{eq:l8-F-contact}がそれらの位置を一致させるので、結果はまさに四スカラー接触頂点である。

\begin{figure}[H]
\centering
\begin{tikzpicture}[scale=0.95]
  \coordinate (vfour) at (0,0);
  \foreach \angle in {35,145,215,325}
    \draw[thick,dashed] (vfour) -- (\angle:1.9);
  \fill (vfour) circle (1.5pt);
  \node[below=8pt of vfour] {\(-|g|^2|\phi|^4\)};
\end{tikzpicture}
\caption{補助場が消去された後、2つの局所的な\(F\)頂点が1つの四スカラー接触相互作用になる。}
\end{figure}

超対称性は、ボソンとフェルミオンを同じ理論に置く以上のことをした。それはそれらの質量を一致させ、湯川結合を\(g\)に結びつけ、四次スカラー結合を\(|g|^2\)に固定した。これらの関係は質量、散乱振幅、崩壊を通じて検証されうる；それらは任意の数値的な偶然ではない。

\section{対になったループ効果}

これらの関係は、相互作用が粒子の質量を補正するときに最も有用になる。スカラー自己エネルギーは、四点頂点からのボソンループと、二つの湯川頂点からのフェルミオンループを受け取る。それぞれが同じ結合定数の二乗を運ぶ。閉じたフェルミオンループは、フェルミ統計によって要求される相対的な負符号を持つ。

\begin{figure}[H]
\centering
\begin{tikzpicture}[>=Latex,scale=0.9]
  \begin{scope}[xshift=-3.3cm]
    \draw[thick,dashed] (-2,0)--(0,0)--(2,0);
    \draw[thick,dashed] (0,0) .. controls (-0.9,1.7) and (0.9,1.7) .. (0,0);
    \fill (0,0) circle (1.4pt);
    \node at (0,-0.55) {bosonic loop};
  \end{scope}
  \begin{scope}[xshift=3.3cm]
    \draw[thick,dashed] (-2,0)--(-0.8,0);
    \draw[thick,dashed] (0.8,0)--(2,0);
    \draw[->,thick] (-0.8,0) arc[start angle=180,end angle=0,x radius=0.8,y radius=0.8];
    \draw[->,thick] (0.8,0) arc[start angle=0,end angle=-180,x radius=0.8,y radius=0.8];
    \fill (-0.8,0) circle (1.4pt);
    \fill (0.8,0) circle (1.4pt);
    \node at (0,-1.35) {fermionic loop};
  \end{scope}
\end{tikzpicture}
\caption{講義で対にされたスカラー自己エネルギーの機構。超対称性が相対的な結合定数を固定し、閉じたフェルミオンループが反対の符号を与える。}
\end{figure}

1ループでは、主要な紫外に敏感な部分が、完全なボソン的寄与とフェルミオン的寄与の間で相殺する。同じ対称性が、Ward恒等式と非繰り込み性によって高次を組織する。この主張は対称性に関連した和についてのものであり、すべてのグラフが文字通り1対1で相殺するということではない。波動関数の繰り込み、有限の効果、その他の対称性を保つ補正は残りうる。これが超空間で理論を構築する実際的な力である。すなわち、質量の等しさと結合定数の間の関係が、個々のファインマン積分を評価する前に強制される。\editorialnote{黒板は相殺をより広く述べている。本文は標準的な限定を与えている。}

\section{次に来るもの}

ここで構築された模型は最も単純な相互作用するカイラル理論であり、一般にWess--Zumino模型と呼ばれる。それは一例であり、最も一般的な超対称場の理論ではないが、パターンは持続する。すなわち、場を超多重項で表し、不変な超空間積分を書き、グラスマン展開に成分の関係を強制させる。

厳密な超対称性は観測された粒子スペクトルには見えないので、次の問題は、有用であるのに十分な紫外制御を保ちながら、どのように対称性を破ることができるかである。もう一つの約束された方向は大統一である。どのように
\begin{equation}
SU(3)\times SU(2)\times U(1)
\end{equation}
が\(SU(5)\)のようなより大きな群に収まるのか、そして超対称な粒子内容がその統一の話をどのように変えるのか。これらが残りの講義が向かう問いである。
